\par\SourceLeaf{127}
Daraus wird deutlich, dass sich auch die Verbindlichkeit II. mit ziemlicher Zuverlässigkeit in Geld bewerten lässt und dass zur Durchführung dieser Rechnungen für I. und II. Geburtsjahr und Geburtstag folgender Personen beschafft werden müssen:
\begin{quote}
der einzelnen jetzt lebenden Witwen,\par
ihrer Kinder unter 20 Jahren, soweit solche vorhanden sind, wie bei Frau Hofr. M., Frau G. J. R. M. und Frau Prof. H.,\par
der jetzt verheirateten Mitglieder,\par
ihrer Ehefrauen.
\end{quote}
\noindent\hspace*{\parindent}Ad III. Am bedenklichsten muss aber das Vorhaben erscheinen, den gegenwärtigen Geldwert der Verbindlichkeit der Kasse gegenüber den künftigen Teilnehmern \textit{in saecula saeculorum} in Zahlen auszudrücken, selbstverständlich nach den jetzigen Statuten und bei der jetzigen Höhe der Pensionen und Beiträge. Und doch ist es \textit{notwendig}, dass man sich davon wenigstens eine ungefähre Vorstellung machen kann. Denn es geht ja gerade darum, \textit{dass die Stabilität} nicht einer Witwenkasse, deren Einrichtungen demnächst je nach den Umständen geändert werden sollen, sondern \textit{unserer Witwenkasse nach ihren jetzigen Einrichtungen beurteilt werden soll.} Dabei wird man natürlich nicht vergessen, dass die Rechnung von gewissen Elementen abhängig bleibt, die sich teils schon jetzt nur annähernd abschätzen lassen, teils im Lauf der Zeit sehr erhebliche Veränderungen erfahren können. Von diesen Elementen nenne ich zwei: die Höhe des Zinssatzes und die Zahl der im Jahresdurchschnitt neu hinzutretenden Mitglieder.

Bei einer Anstalt, die nur zu einem sehr kleinen Teil auf laufenden Beiträgen, hauptsächlich aber auf Kapitalerträgen beruht und beruhen soll, steht die Höhe des Zinssatzes offenbar in einem genauen (umgekehrten) Verhältnis zur Höhe des erforderlichen Kapitals. Wenn beispielsweise zu einem bestimmten Zeitpunkt die Schenkung eines Kapitals von 70000 Thl. gerade ausreichte, um bei einem Zinssatz von 4 p.c. eine Gesellschaft zu erhalten, die im Durchschnitt jedes Jahr gleich viele neue Mitglieder aufnimmt, würde das Sinken des Zinssatzes auf \(3\frac12\) p.c. eine Erhöhung des Kapitals auf 80000 Thl. erfordern. Ich halte diesen Umstand gerade \textit{im Hinblick auf die Schwierigkeit der Beurteilung} für den unerheblichsten. Die Beurteilung kann nicht mehr tun, als die Größe dieses Einflusses klar darzustellen. Daraus folgt von selbst, dass bereits dafür gesorgt werden muss, dass das Kapital im Jahresdurchschnitt einen angemessenen Zuwachs erhält, um dem im Lauf der Zeit allmählich, aber unfehlbar eintretenden Sinken des Zinssatzes zu begegnen.

Ebenso einleuchtend ist, dass die Höhe des erforderlichen Kapitals der durchschnittlichen Zahl der jährlich neu beitretenden Mitglieder (\textit{ceteris paribus}) genau proportional sein wird. Zunächst können wir nur unsere eigenen Erfahrungen zugrunde legen, die seit 100 Jahren vorliegen und bei denen natürlich die neueren und neuesten Zeiten unser Urteil vorzugsweise leiten müssen. Seine Magnifizenz bemerkt mit Recht, dass die aus gestiegenen wissenschaftlichen Bedürfnissen und Anforderungen hervorgegangene Vergrößerung der Zahl der Professoren einen wesentlichen Einfluss auf das Bestehen solcher Professorenwitwenkassen haben muss, die hauptsächlich auf Kapital gegründet sind. Es ist daher sehr wohl möglich, dass die Göttinger Ergebnisse etwa aus den letzten 30 oder 40 Jahren keinen ganz richtigen Maßstab für die Zukunft, besonders für die Zukunft späterer Jahrhunderte, bilden können. Doch diese Ungewissheit liegt in der Natur der Veränderlichkeit aller menschlichen Dinge; ihre Folgen treten \textit{allmählich} hervor, und man begegnet ihnen nur mit einer Wachsamkeit, die niemals einschläft. In unserem Fall berechnet man also das Kapital, das zur dauerhaften Erfüllung der Verbindlichkeit III. erforderlich ist, nach unseren besten gegenwärtigen Kenntnissen, vergisst nicht, dass gerade wegen jener Ungewissheit ein etwas größeres Kapital vorhanden sein muss, \textit{wiederholt} die Rechnung fortlaufend in bestimmten, nicht allzu großen
\SourceLeaf{128}
Zeitabständen, etwa alle 5 oder 10 Jahre, zieht dabei stets die neu hinzugekommenen Erfahrungen heran und betrachtet den Überschuss erst dann als teilweise verfügbar, wenn er sich wirklich \textit{vergrößert} hat.

Aber auch abgesehen von diesen beiden Umständen, oder mit anderen Worten: Selbst wenn man einen bestimmten Zinssatz und eine bestimmte durchschnittliche Zahl der jährlich neu beitretenden Mitglieder zugrunde legt, scheint die Schwierigkeit der Abschätzung doch fast unüberwindlich. Denn die verschiedensten Altersverhältnisse zwischen den Ehegatten kommen vor, von der Wiederverheiratung Verwitweter gar nicht erst zu reden. In dieser Hinsicht scheint der Gutachter also ungefähr auf derselben Stufe zu stehen wie jemand, der im Voraus den Plan einer der vielen Witwenkassen hätte prüfen sollen, die ohne strenge Berücksichtigung des Lebensalters der eintretenden Ehepaare errichtet wurden und fast alle zugrunde gegangen sind. (Wenn Herr Universitätsrat K. glaubt, es werde auch bei all diesen Kassen an \textit{Berechnungen} nicht gefehlt haben, hat er zweifellos recht; wenn er daraus aber auf die Grundlosigkeit der \textit{Wahrscheinlichkeitsrechnung} schließen will, hat er unrecht. Gewiss gibt es viele Wörter, mit denen verschiedene Personen verschiedene Bedeutungen verbinden, und ebenso solche, die wissenschaftlich eine sehr bestimmte Bedeutung haben, unter denen man im Alltag aber oft sehr Verschiedenes zusammenfasst. So verhält es sich mit dem Ausdruck Wahrscheinlichkeitsrechnung. Im strengen Sinn verstanden, kann von ihrer Anwendung in allen Fällen, in denen die nötigen Grundlagen fehlen, gar nicht die Rede sein. Bei der Einrichtung all der gescheiterten Witwenkassen wurde von der strengen Wahrscheinlichkeitsrechnung \textit{überhaupt kein} Gebrauch gemacht, sondern nur von vagen Einschätzungen. Das spreche ich hier nur als Tatsache aus, nicht als Vorwurf, denn eine Grundlage in der Wahrscheinlichkeitsrechnung war schon deshalb tatsächlich unmöglich, weil alle nötigen Voraussetzungen dafür fehlten.) Im vorliegenden Fall ist es zwar unmöglich, aus den einzelnen Elementen nach der Wahrscheinlichkeitsrechnung ein Endergebnis zu ermitteln, gerade weil diese Elemente fehlen. Wohl aber bietet die hundertjährige Erfahrung der Kasse selbst einen reichen Schatz als Grundlage, wenn man sie auf die richtige Weise nutzt. Diese Erfahrungen müssen daher erst gesammelt und geordnet werden. Ich setze voraus, dass ein Buch angelegt wird, in dem seit der ersten Gründung der Gesellschaft sämtliche Mitglieder ohne Ausnahme in der zeitlichen Reihenfolge ihres Eintritts verzeichnet werden, zusammen mit allen Angaben, die für den fraglichen Zweck erheblich sind. Diese Erfahrungen würden allerdings noch viel ergiebigeres Material bieten, wenn für sämtliche beteiligten Personen auch Geburtsjahr und Geburtstag aufgezeichnet wären: für den eintretenden Professor, für seine Frau, falls er schon verheiratet ist oder wenn und sooft er sich nach dem Eintritt verheiratet, und schließlich für die minderjährigen Kinder, die beim Tod des Mitglieds leben. All diese Angaben fehlen aber und ließen sich jetzt in den meisten Fällen nur noch für das Mitglied selbst ergänzen; aus diesen vereinzelten Angaben ließe sich jedoch wenig oder gar kein Nutzen ziehen. Dennoch bleibt das, was sich zweifellos noch jetzt zusammentragen lässt, höchst wertvoll. Ich meine für jedes einzelne Mitglied:

1. \textit{Terminus a quo} und \textit{ad quem} der geleisteten Beiträge.

2. \textit{Terminus a quo} und \textit{ad quem} der bezogenen Witwenpension, in den Fällen, in denen eine solche eintrat.

3. \textit{Terminus a quo} und \textit{ad quem} der bezogenen Waisenpension, wenn nach dem Tod des Vaters oder der Mutter noch minderjährige Kinder vorhanden waren.

Die \textit{Höhe} der Geldbeträge, die von den Mitgliedern beigetragen und von den Witwen und Waisen bezogen wurden, braucht aber gar nicht mit herausgeschrieben zu werden.

Dies wäre also die \textit{dritte Voraussetzung}, deren Beschaffung vor Beginn aller Berechnungen unerlässlich ist. Ich kenne die Einrichtung des Archivs der Witwenkasse überhaupt nicht und weiß daher nicht, ob vielleicht bereits eigens angelegte Bücher vorhanden sind, denen sich dieses Material mit Sicherheit, Vollständigkeit
\SourceLeaf{129}
und Leichtigkeit entnehmen lässt. Jedenfalls könnten aber die zweifellos aufbewahrten 105 Jahresrechnungen dazu dienen. Erst nachdem ich selbst die vorhandenen Unterlagen näher eingesehen habe, könnte ich mich dazu äußern, ob und in welchem Maß ich meine eigene Mithilfe bei dieser Auswertung zusagen kann.

Für jeden einzelnen Teilnehmer, der in diesem Buch aufzuführen ist, lässt sich aus den nach Rubriken geordneten Daten berechnen, wie viel er der Gesellschaft und wie viel diese seinen Hinterbliebenen bei den gegenwärtigen Sätzen geleistet haben und welchen auf den Zeitpunkt seines Eintritts zurückgerechneten Geldwert dies bei einem bestimmten Zinssatz gehabt hätte. Natürlich ist dies bei den einzelnen Personen sehr verschieden, bei einigen positiv, bei anderen negativ. Nach den ewigen Gesetzen ist aber der \textit{Mittelwert} aus einer großen Anzahl ein Element, das wiederum als Mittelwert für die Zukunft zugrunde gelegt werden kann, wenn kein Grund besteht, wesentliche Änderungen der allgemeinen Verhältnisse vorauszusetzen. Schon diese Tabelle selbst kann darüber einen lehrreichen Hinweis geben, wenn man das Ganze gruppiert und etwa den Mittelwert der ersten Hälfte mit dem der zweiten oder das erste, zweite und dritte Drittel miteinander vergleicht. So wird sich herausstellen, wie groß der Geldwert der Verbindlichkeit der Gesellschaft ist, die ihr durch den Eintritt eines neuen Mitglieds \textit{durchschnittlich} zuwächst. Wenn man zugleich nach dem oben schon Bemerkten eine plausible Annahme über die durchschnittliche Zahl der jährlich Hinzutretenden getroffen hat, lässt sich für einen bestimmten Zinssatz berechnen, wie groß das Kapital sein muss, dessen Zinsertrag diese für immer fortlaufende Verbindlichkeit III. decken kann. Alle drei Kapitalbeträge für I., II. und III. zusammengenommen und mit dem tatsächlichen Vermögen der Gesellschaft verglichen, werden dann so genau, wie es die Natur des Gegenstands zulässt, zeigen, ob bei der jetzigen Einrichtung ihre Stabilität \textit{mehr} als gesichert ist oder nur eben ausreicht oder ob sich daraus ihre Instabilität ergibt und deshalb entschieden früher oder später ihr Untergang erwartet werden muss. Dementsprechend wären dann in den verschiedenen Fällen die geeigneten Maßnahmen zu erwägen und einzuleiten.

Dies sind die Grundzüge des Plans, nach dem meiner Meinung nach eine gründliche Prüfung und die Aufstellung einer Bilanz ausgeführt werden könnten und müssten. Es ist eine bedeutende Arbeit, die ich aber gern übernehmen werde, wenn es gewünscht wird. Dass ich diese Skizze hier nur ungefeilt und flüchtig niedergeschrieben vorlege, wird man mit der Kürze der Zeit entschuldigen. Wird eine solche Arbeit jetzt ausgeführt, bleibt es jedenfalls, wie schon oben bemerkt, dringend wünschenswert, dass künftig in bestimmten Zeitabständen immer wieder eine neue Bilanz gezogen wird. Das würde dann viel leichter als beim ersten Mal, wenn ein solches Buch, wie ich es oben erwähnt habe, wenigstens von jetzt an mit allen \textit{Datumsrubriken} angelegt, regelmäßig vervollständigt und fortgeführt würde.

Ich will nun auch noch mein Votum zu ein paar anderen Punkten beifügen, die in den anderen Abstimmungen berührt worden sind.

Ich bin nicht dafür, die Bestimmung der Statuten aufzuheben, die unbesoldete Professoren ausschließt. Die Universität als Körperschaft müsste im Hinblick auf die Witwenkasse stets dringend wünschen, dass Fälle, in denen einem bei Schule oder Kirche Angestellten der Professortitel verliehen wird, sehr selten bleiben. Von einem solchen Fall abgesehen, wird jemand, der weder Besoldung noch Vermögen hat, nicht leicht so unbesonnen sein, eine Ehe einzugehen, sodass die ganze Bestimmung überhaupt selten, vielleicht nie, eine Wirkung haben wird. Möglicherweise könnte aber ein unbesoldeter Professor, der sich dem Tod nahe fühlt und eine Braut hat, die er sonst vor Erlangung einer Besoldung gewiss noch nicht geheiratet hätte, falls ihm der Eintritt in die Witwenkasse offenstünde, dadurch versucht werden, der Witwenkasse durch eine rasche Trauung eine Last aufzubürden.

Solange \textit{bona fide} gehandelt wird, müssen vielmehr die unbesoldeten Professoren jene Klausel als
\SourceLeaf{130}
eine gerechte, zu ihrem Vorteil gereichende Bestimmung betrachten. Sie erspart ihnen die Alternative, entweder schon während der Zeit, in der sie nichts von der Universität erhalten, zur Witwenkasse beitragen zu müssen oder später, wenn sie eine Besoldung erhalten, noch für die ganze Zeit ihres unbesoldeten Professorenstandes doppelt nachzahlen zu müssen.

Meine zweite Bemerkung betrifft den Zinssatz. Hier kann ich dem, was im P. M. des U. R. O. gesagt ist, nicht ganz zustimmen. Vielmehr erscheint mir die Rechnung, nach der der jetzige Zinssatz mit \(4\tfrac1{10}\) Prozent ermittelt wurde, zum Teil als illusorisch. Ich erläutere das an einem Beispiel. Die österreichischen \(4\tfrac12\)-Prozent-Papiere stehen nach dem heutigen Kurszettel bei \(103\tfrac58\). Beim Kauf von einem Bankier wird man alles eingerechnet gewiss \textit{mehr als} 104 tatsächlich zahlen müssen; ich will aber nur bei 104 bleiben. Man erhält also für sein eingezahltes Geld tatsächlich nur \(4\tfrac{17}{32}\) oder nicht ganz \(4\tfrac13\) Prozent Zinsen. Es dauert somit mindestens 12 Jahre, bis man überhaupt sagen kann, dass man tatsächlich 4 Prozent Zinsen erhalten hat. Nun werden aber von diesen Papieren jedes Jahr \textit{sehr große} Summen ausgelost und zum Nennwert zurückgezahlt. Erfolgt die Auslosung schon nach 2 Jahren, hat man tatsächlich nur insgesamt \(4\tfrac23\) Prozent oder für ein Jahr \(2\tfrac13\) Prozent Zinsen erhalten, ohne die Kosten zu rechnen, die mit jeder Einziehung verbunden sind. Für den Besitzer eines solchen Papiers ist es auch immer ein gefährlicher Umstand, dass er einen sehr bedeutenden Verlust erleiden kann, wenn die ihn betreffende Auslosung ihm nicht bekannt wird und er dadurch die rechtzeitige Einziehung versäumt. Für die Witwenkasse wird der Bankier, von dem die Papiere gekauft wurden, wohl immer die nötige Wachsamkeit ausüben, weil ihm selbst jedes anfallende Geschäft einen Gewinn einbringt. Aber die eigentliche Verantwortung für jeden durch ein mögliches Übersehen entstandenen Verlust wird er doch kaum übernehmen. In dieser Hinsicht will ich daher nicht unterlassen, hiermit mitzuteilen, dass bei der Auslosung von \textit{anderthalb Millionen Gulden} der fraglichen Papiere, die heute vor acht Tagen in Wien stattgefunden hat, auch eine der Obligationen der Witwenkasse gezogen wurde: die auf Seite IX. der Rechnung unter Nr. 52 aufgeführte Litr. P Nr. 15472. Dass ich diese Mitteilung machen kann, verdanke ich nur dem zufälligen Umstand, dass ich heute, wo gerade diese Rechnung in meinen Händen ist, die Nachricht von der erfolgten Auslosung in einer Zeitung fand und mir daher die Bezeichnung der ausgelosten Nummer notierte, um sie zu Hause mit der Kapitalliste der Witwenkasse vergleichen zu können. Mit dieser Mitteilung will ich nun diese lange Darlegung beschließen.

\noindent 9. Januar 1845.\hfill\textsc{Gauss.}

